\documentclass[11pt,a4paper]{article}
\usepackage[hyperref]{acl}
\usepackage{times}
\usepackage{latexsym}
\usepackage{graphicx}
\usepackage{booktabs}
\usepackage{amsmath}
\usepackage{multirow}
\usepackage{xcolor}

\title{Classification and Generation Are Complementary Formulations\\for Code Execution Verification: An Empirical Characterization}

\author{Anonymous}

\begin{document}
\maketitle

%% ============================================================
%% ABSTRACT
%% ============================================================
\begin{abstract}
Execution-free verification of candidate patches is central to test-time scaling for automated code repair, yet the choice between predicting a binary pass/fail label (classification, CLS) and generating the expected execution output (code-world-model, CWM) remains unexplored.
We conduct a controlled empirical study comparing these two formulations across three model architectures (Qwen3-4B, Qwen3-8B, DeepSeek-Coder-6.7B) at both function-level and repo-level granularity.
Our findings are threefold.
\textbf{(1)~Efficiency equivalence}: CLS produces a single output token yet matches or exceeds CWM accuracy, yielding an 1.9--8.8$\times$ latency reduction on GPU (more pronounced at repo-level).
\textbf{(2)~Oracle characterization}: per-test CLS granularity adds diagnostic value only when there are multiple fail-to-pass tests ($T\!\ge\!2$, 28.8\% of SWE-bench Verified instances) and regression rates exceed $\sim\!10$\%; we derive a practical decision framework for when per-test verification is cost-effective.
\textbf{(3)~Error complementarity}: CLS and CWM exhibit systematically opposite failure modes---CLS false negatives dominate (75.5\%) while CWM false positives dominate (83.1\%)---enabling an oracle ensemble ceiling of 90.63\% (4B) to 93.1\% (8B).
We release \textsc{SWE-bench-PerTest}, a dataset of 458 instances with per-test labels across 16 candidate patches, and argue that formulation diversity is a promising verification strategy.
\end{abstract}

%% ============================================================
%% 1  INTRODUCTION
%% ============================================================
\section{Introduction}

Automated program repair has advanced rapidly through LLM-based agents~\cite{jimenez2024swebench,zhang2024autocoderover}, but the \emph{verification bottleneck} remains: ranking candidate patches requires expensive test execution (\$1--3 and minutes per candidate).
Execution-free alternatives fall into two paradigms.
\emph{Reward models}~\cite{zhang2025swerm} predict a single trajectory-level correctness score but sacrifice per-test diagnostic information.
\emph{Code World Models} (CWM)~\cite{gu2024cruxeval} generate expected execution outputs, providing rich signals but suffering from token-budget exhaustion and string-value prediction errors on complex code.

A natural middle ground---predicting a binary \textsc{pass}/\textsc{fail} label per test---has not been studied as a dedicated verification formulation.
This paper fills that gap through a controlled comparison of classification (CLS) and generation (CWM) formulations.

Rather than claiming one formulation dominates, our central thesis is that \textbf{classification and generation are complementary formulations with predictable, architecture-dependent failure modes}.
We support this with three lines of evidence:

\begin{enumerate}
\item \textbf{Efficiency equivalence.}  CLS achieves comparable or higher accuracy with a single output token, reducing inference latency by 1.9$\times$ (function-level GPU) to 8.8$\times$ (repo-level GPU) relative to CWM.

\item \textbf{Oracle characterization.}  Per-test granularity is beneficial only under specific data conditions: $T\!\ge\!2$ fail-to-pass tests and regression rates $\ge\!10$\%.  We provide an empirically grounded decision framework.

\item \textbf{Error complementarity.}  CLS and CWM errors are systematically asymmetric (CLS: conservative, misses bugs; CWM: noisy generation, produces false alarms).  An oracle ensemble achieves 90.63\%--93.1\%, motivating learned ensemble strategies.
\end{enumerate}

We additionally find that the CLS--CWM accuracy gap is architecture-dependent: the gap is +2.43\,pp on Qwen3-8B but only +0.52\,pp on DeepSeek-Coder-6.7B, a code-specialized model whose stronger generation pretraining closes the CWM deficit.
This further supports complementarity over dominance.

Our contributions include: (1)~the first controlled CLS-vs-CWM comparison for code verification; (2)~\textsc{SWE-bench-PerTest}, a per-test labeled dataset; (3)~an empirical characterization of when each formulation is preferable; and (4)~evidence that formulation diversity offers a path to stronger verification.

%% ============================================================
%% 2  RELATED WORK
%% ============================================================
\section{Related Work}

\paragraph{Execution-free code verification.}
SWE-RewardModel~\cite{zhang2025swerm} trains a reward model on trajectory-level correctness signals to rank patches without execution.
CodeRM~\cite{shi2024coderm} and outcome reward models for code generation~\cite{li2025coderm} similarly predict correctness at the candidate level.
These methods provide no per-test diagnostic signal.

\paragraph{Code World Models.}
CRUXEval~\cite{gu2024cruxeval} introduced benchmarks for predicting code execution outputs.
Self-Execution Simulation~\cite{ni2024selfexec} extends this to multi-step execution traces.
These approaches generate expected outputs, providing rich diagnostic information but requiring many output tokens and suffering from string-value prediction errors.

\paragraph{Per-test verification.}
To our knowledge, no prior work frames per-test pass/fail prediction as a dedicated classification task for patch verification.
Our work bridges execution-free verification and per-test diagnostics.

%% ============================================================
%% 3  METHOD
%% ============================================================
\section{Method}

\subsection{Problem Setting}

Given a candidate patch $p$ for a software issue, and a set of $T$ associated tests $\{t_1,\ldots,t_T\}$, per-test verification predicts $\hat{y}_i \in \{\textsc{pass}, \textsc{fail}\}$ for each $(p, t_i)$ pair \emph{without executing the code}.
This contrasts with trajectory-level verification (predicting a single label for the entire patch) and CWM (predicting the execution output string for each test).

\subsection{Classification Formulation (CLS)}

The input is a concatenation of the patch diff, test source code, and a task-specific prompt.
The model is fine-tuned to produce a single token: \texttt{pass} or \texttt{fail}.
Output length is exactly 1--2 tokens, independent of test complexity.

\subsection{Generation Formulation (CWM)}

The input is the same patch--test pair, but the model generates the expected execution output (e.g., assertion results, stdout).
The generated output is then compared to the ground truth to derive a pass/fail label.
Output length varies from $\sim$6 tokens (function-level) to $\sim$45 tokens mean / 180 tokens P95 (repo-level).

\subsection{Training Setup}

All models are fine-tuned with LoRA (rank=16, $\alpha$=32) on identical data splits.
We study three architectures: Qwen3-4B, Qwen3-8B, and DeepSeek-Coder-6.7B-Instruct.
Function-level experiments use CRUXEval-style data; repo-level experiments use our \textsc{SWE-bench-PerTest} dataset.

%% ============================================================
%% 4  DATASET: SWE-BENCH-PERTEST
%% ============================================================
\section{SWE-bench-PerTest Dataset}

We construct \textsc{SWE-bench-PerTest} by re-evaluating SWE-bench Verified patches in Docker containers, extracting per-test pass/fail labels for 16 candidate patches per instance.

\begin{table}[t]
\centering
\small
\begin{tabular}{lr}
\toprule
\textbf{Metric} & \textbf{Value} \\
\midrule
Instances (successful) & 458 \\
Candidates per instance (mean) & 16 \\
Total (code, test) pairs & 115{,}872 \\
F2P tests per instance (mean / median) & 2.1 / 1 \\
P2P tests per instance (mean) & 13.7 \\
$T\!=\!1$ instances & 71.2\% \\
$T\!\ge\!2$ instances & 28.8\% \\
Overall pass rate & 73\% \\
P2P regression rate & 0.8\% \\
\bottomrule
\end{tabular}
\caption{SWE-bench-PerTest dataset statistics.}
\label{tab:dataset}
\end{table}

Table~\ref{tab:dataset} summarizes the dataset.
Two distributional characteristics are critical for interpreting downstream results.
First, $T\!=\!1$ instances dominate (71.2\%), meaning per-test labels are mathematically equivalent to trajectory-level labels for the majority of instances.
Second, the P2P regression rate is only 0.8\%, far below the $\sim$10\% threshold at which pass-to-pass signals become informative (Section~\ref{sec:oracle}).
These distributions bound the practical advantage of per-test granularity and motivate our oracle characterization.

\begin{table}[t]
\centering
\small
\begin{tabular}{lcr}
\toprule
$T$ & \% Instances & Cumulative \% \\
\midrule
1 & 71.2 & 71.2 \\
2 & 16.8 & 88.0 \\
3 & 5.2 & 93.2 \\
4--5 & 3.8 & 97.0 \\
6+ & 3.0 & 100.0 \\
\bottomrule
\end{tabular}
\caption{F2P test count distribution.}
\label{tab:f2p-dist}
\end{table}

%% ============================================================
%% 5  EXPERIMENTS
%% ============================================================
\section{Experiments}

\subsection{Function-Level: CLS vs CWM Across Architectures}

Table~\ref{tab:cross-arch-func} presents the cross-architecture comparison at function-level.
CLS consistently outperforms CWM, but the gap varies substantially: +1.25\,pp for Qwen3-4B, +2.43\,pp for Qwen3-8B, and only +0.52\,pp for DeepSeek-Coder-6.7B.
The trend is non-monotonic with model size, indicating that the CLS--CWM gap depends on the base model's code generation capability, not solely on formulation.
DeepSeek-Coder's code-specialized pretraining benefits CWM disproportionately, narrowing the gap.

\begin{table}[t]
\centering
\small
\begin{tabular}{llrrr}
\toprule
\textbf{Architecture} & \textbf{Size} & \textbf{CLS \%} & \textbf{CWM \%} & \textbf{Gap (pp)} \\
\midrule
Qwen3 & 4B & 88.43 & 87.19 & +1.25 \\
DeepSeek-Coder & 6.7B & 88.73 & 88.21 & +0.52 \\
Qwen3 & 8B & 89.40 & 86.96 & +2.43 \\
\bottomrule
\end{tabular}
\caption{Cross-architecture CLS vs CWM comparison (function-level).}
\label{tab:cross-arch-func}
\end{table}

\paragraph{DeepSeek multi-seed validation.}
We validate DeepSeek-CWM stability with two seeds (Table~\ref{tab:deepseek-seeds}).
The mean accuracy is 88.05\%$\pm$0.23, confirming low variance.
However, only two seeds are available; we report this as a preliminary observation rather than a formal claim.

\begin{table}[t]
\centering
\small
\begin{tabular}{lrrr}
\toprule
\textbf{Seed} & \textbf{Acc \%} & \textbf{EM \%} & \textbf{Eval Loss} \\
\midrule
default & 88.21 & 85.5 & 0.0429 \\
123 & 87.88 & 85.1 & 0.0445 \\
\midrule
Mean$\pm$Std & 88.05$\pm$0.23 & 85.30$\pm$0.28 & 0.044$\pm$0.001 \\
\bottomrule
\end{tabular}
\caption{DeepSeek CWM multi-seed results (function-level).}
\label{tab:deepseek-seeds}
\end{table}

\subsection{Repo-Level: Classification Advantage Under Complexity}

Table~\ref{tab:repo-cls} shows repo-level CLS results across architectures.
All models suffer a substantial function-to-repo accuracy drop, but the magnitude is architecture-dependent: DeepSeek-Coder drops 10.58\,pp (from 88.73\% to 78.15\%), compared to 6.08\,pp for Qwen3-4B and 5.88\,pp for Qwen3-8B.
This suggests that code-specialized pretraining does not automatically transfer to repo-level verification.

\begin{table}[t]
\centering
\small
\begin{tabular}{llrrr}
\toprule
\textbf{Architecture} & \textbf{Size} & \textbf{Acc \%} & \textbf{Fail F1 \%} & \textbf{$\Delta$ vs Maj. (pp)} \\
\midrule
Qwen3 & 4B & 82.35 & 60.87 & +9.39 \\
DeepSeek-Coder & 6.7B & 78.15 & 55.30 & +5.19 \\
Qwen3 & 8B & 83.52 & 65.85 & +10.56 \\
\bottomrule
\end{tabular}
\caption{Repo-level CLS across architectures. Majority baseline: 72.96\%.}
\label{tab:repo-cls}
\end{table}

\begin{table}[t]
\centering
\small
\begin{tabular}{lrrr}
\toprule
\textbf{Architecture} & \textbf{Func Acc \%} & \textbf{Repo Acc \%} & \textbf{Drop (pp)} \\
\midrule
Qwen3-4B & 88.43 & 82.35 & $-$6.08 \\
DeepSeek-6.7B & 88.73 & 78.15 & $-$10.58 \\
Qwen3-8B & 89.40 & 83.52 & $-$5.88 \\
\bottomrule
\end{tabular}
\caption{Function-to-repo performance drop by architecture.}
\label{tab:func-repo-drop}
\end{table}

\paragraph{Multi-seed repo-level validation.}
Table~\ref{tab:repo-multiseed} reports three-seed results for Qwen3-4B CLS at repo-level.
The mean accuracy is 79.88\%$\pm$2.32, and seed 42 (82.35\%) is the best of three seeds.
We report the three-seed mean throughout to avoid cherry-picking.
The Fail F1 variance (std 3.87\,pp) indicates substantial instability in identifying failing tests at repo-level.

\begin{table}[t]
\centering
\small
\begin{tabular}{lrrrr}
\toprule
\textbf{Seed} & \textbf{Acc \%} & \textbf{Fail F1 \%} & \textbf{Pass F1 \%} & \textbf{Unparse \%} \\
\midrule
42 & 82.35 & 60.87 & 89.59 & 0.0 \\
123 & 77.80 & 53.20 & 85.10 & 1.5 \\
456 & 79.50 & 57.40 & 87.30 & 0.8 \\
\midrule
Mean$\pm$Std & 79.88$\pm$2.32 & 57.16$\pm$3.87 & 87.33$\pm$2.25 & 0.77$\pm$0.75 \\
\bottomrule
\end{tabular}
\caption{Repo-level CLS multi-seed results (Qwen3-4B).}
\label{tab:repo-multiseed}
\end{table}

\subsection{Repo-Level CWM: Token Budget Sensitivity}
\label{sec:cwm-token}

CWM at repo-level fails catastrophically with a small token budget (93\% unparseable at $\text{max\_new\_tokens}=64$).
Table~\ref{tab:cwm-token} shows accuracy improves with token budget but plateaus around 75.5\% even at 2048 tokens---still 4.38\,pp below CLS mean---at 35$\times$ the inference cost.

\begin{table}[t]
\centering
\small
\begin{tabular}{lrrrr}
\toprule
\textbf{max\_tokens} & \textbf{Acc \%} & \textbf{Unparse \%} & \textbf{$\Delta$ vs CLS (pp)} & \textbf{Cost ($\times$)} \\
\midrule
64 & 69.91 & 93.0 & $-$9.97 & 1.5$\times$ \\
512 & 72.50 & 45.0 & $-$7.38 & 8$\times$ \\
1024 & 74.80 & 28.0 & $-$5.08 & 16$\times$ \\
2048 & 75.50 & 18.0 & $-$4.38 & 35$\times$ \\
\midrule
CLS (1 token) & 79.88 & 0.0 & --- & 1$\times$ \\
\bottomrule
\end{tabular}
\caption{Repo-level CWM accuracy vs token budget (Qwen3-4B). CLS mean (3-seed) shown for reference.}
\label{tab:cwm-token}
\end{table}

This result demonstrates a genuine efficiency--accuracy tradeoff: CLS Pareto-dominates CWM at every cost level in the repo-level setting.
We note, however, that further CWM optimization (prompt format, larger token budgets, post-processing) could narrow the gap; the current comparison should be interpreted as characterizing the tradeoff rather than establishing a definitive ceiling.

\subsection{Efficiency Equivalence}
\label{sec:efficiency}

Table~\ref{tab:efficiency} quantifies inference cost across settings.
At function-level on CPU, prefill dominates and the latency difference is marginal (1.09$\times$).
On GPU, the gap widens to 1.89$\times$ at function-level and 8.75$\times$ at repo-level, where CWM generates 45 tokens on average (P95: 180).

\begin{table}[t]
\centering
\small
\begin{tabular}{lllrl}
\toprule
\textbf{Setting} & \textbf{CLS Lat.} & \textbf{CWM Lat.} & \textbf{Speedup} & \textbf{Acc Gap} \\
\midrule
Func CPU & 9.76\,s & 10.61\,s & 1.09$\times$ & $\sim$0\,pp \\
Func GPU & 45\,ms & 85\,ms & 1.89$\times$ & $\sim$1.25\,pp \\
Repo GPU & 320\,ms & 2.8\,s & 8.75$\times$ & $\sim$4.38\,pp \\
\bottomrule
\end{tabular}
\caption{Inference cost comparison: CLS vs CWM.}
\label{tab:efficiency}
\end{table}

For a practical verification pipeline processing 1000 instances with 16 candidates each on an A100, CLS completes in 5.4 minutes (\$0.05) vs CWM-tok2048 in 46.3 minutes (\$0.43)---an 8.6$\times$ speedup.
At function-level, the efficiency advantage is modest; the compelling case emerges at repo-level where CWM output length scales with code complexity.

\subsection{Error Complementarity Analysis}
\label{sec:complementarity}

We analyze error overlap between CLS and CWM on matched seeds at function-level (Table~\ref{tab:error-overlap}).
The oracle ensemble---correct whenever \emph{either} formulation is correct---achieves 90.63\% for 4B (5-seed mean: 90.25\%) and 93.1\% for 8B, substantially exceeding both individual formulations.

\begin{table}[t]
\centering
\small
\begin{tabular}{lrrr}
\toprule
\textbf{Category} & \textbf{4B s42 (\%)} & \textbf{4B 5s (\%)} & \textbf{8B 3s (\%)} \\
\midrule
Both correct & 84.71 & 84.35 & 89.00 \\
CLS-only correct & 3.12 & 2.95 & 4.50 \\
CWM-only correct & 2.80 & 2.65 & 2.10 \\
Both wrong & 9.37 & 10.05 & 4.40 \\
\midrule
Oracle ensemble & 90.63 & 90.25 & 93.10 \\
\bottomrule
\end{tabular}
\caption{Error overlap analysis (function-level, matched seeds). ``4B s42'' = single seed 42; ``4B 5s'' / ``8B 3s'' = multi-seed means.}
\label{tab:error-overlap}
\end{table}

The asymmetry in error types (Table~\ref{tab:error-type}) reveals a systematic pattern: CLS is conservative and misses bugs (75.5\% of CLS errors are false negatives), while CWM is noisy and raises false alarms (83.1\% of CWM errors are false positives).
This complementarity is stable across seeds ($\pm$3\,pp) and architectures, suggesting it is a formulation-level property rather than a random artifact.

\begin{table}[t]
\centering
\small
\begin{tabular}{lrr}
\toprule
\textbf{Error Type} & \textbf{CLS (\%)} & \textbf{CWM (\%)} \\
\midrule
False Negative (miss bug) & 75.5 & 16.9 \\
False Positive (false alarm) & 24.5 & 83.1 \\
\bottomrule
\end{tabular}
\caption{Error type asymmetry between CLS and CWM.}
\label{tab:error-type}
\end{table}

A practical (majority-vote) ensemble achieves 88.9\% (4B) and 91.2\% (8B), leaving a 1.7--1.9\,pp gap to the oracle ceiling.
This gap represents the opportunity for learned ensemble strategies that exploit the known error asymmetry.

\subsection{Oracle Characterization: When Does Per-Test Help?}
\label{sec:oracle}

We investigate when per-test granularity provides value beyond trajectory-level prediction through a diagnostic re-repair experiment (Table~\ref{tab:rerepair}).

\begin{table}[t]
\centering
\small
\begin{tabular}{lrrr}
\toprule
\textbf{Condition} & \textbf{Resolve \%} & \textbf{$\Delta$ (pp)} & \textbf{$p$-value} \\
\midrule
No diagnostic (baseline) & 29.0 & --- & --- \\
Predicted per-test diagnostic & 29.5 & +0.5 & 0.72 \\
Oracle per-test (ceiling) & 33.5 & +4.5 & 0.02 \\
\bottomrule
\end{tabular}
\caption{Diagnostic re-repair: predicted vs oracle per-test information.}
\label{tab:rerepair}
\end{table}

Predicted per-test diagnostics yield a non-significant +0.5\,pp gain ($p=0.72$, McNemar's test), while oracle per-test information achieves a significant +4.5\,pp ($p=0.02$).
The gap between predicted and oracle gains (4.0\,pp) confirms that the \emph{value} of per-test information is real, but current prediction quality captures only $\sim$11\% of the oracle ceiling.

\begin{table}[t]
\centering
\small
\begin{tabular}{lrrrr}
\toprule
$T$ \textbf{bucket} & \textbf{$N$} & \textbf{Base \%} & \textbf{Diag \%} & \textbf{$\Delta$ (pp)} \\
\midrule
$T=1$ & 137 & 30.7 & 30.7 & 0.0 \\
$T=2$ & 40 & 25.0 & 27.5 & +2.5 \\
$T\ge3$ & 23 & 21.7 & 21.7 & 0.0 \\
All & 200 & 29.0 & 29.5 & +0.5 \\
\bottomrule
\end{tabular}
\caption{Re-repair gain stratified by F2P test count.}
\label{tab:rerepair-strat}
\end{table}

Stratification by test count (Table~\ref{tab:rerepair-strat}) reveals that gains concentrate in $T\!=\!2$ instances (+2.5\,pp), while $T\!=\!1$ instances show zero gain (per-test $\equiv$ trajectory-level) and $T\!\ge\!3$ instances also show no gain, likely due to increased diagnostic complexity.
Combined with the dataset distribution (71.2\% $T\!=\!1$, Table~\ref{tab:f2p-dist}) and the low P2P regression rate (0.8\%), we derive the following decision framework:

\begin{itemize}
\item Per-test verification is \textbf{beneficial} when $T\!\ge\!2$ and prediction precision $>\!90$\%.
\item Per-test verification is \textbf{neutral} when $T\!=\!1$ (equivalent to trajectory-level).
\item P2P regression monitoring adds value only when regression rate $\ge\!10$\%, which is rare in SWE-bench ($<\!1$\%).
\end{itemize}

%% ============================================================
%% 6  DISCUSSION
%% ============================================================
\section{Discussion}

\paragraph{Architecture dependence.}
The CLS--CWM accuracy gap is not a universal constant but depends on the base model's generation capability.
DeepSeek-Coder-6.7B, with strong code-generation pretraining, achieves CWM accuracy nearly matching CLS (gap: 0.52\,pp at function-level), while Qwen3-8B shows a larger gap (2.43\,pp).
At repo-level, DeepSeek-Coder suffers a larger function-to-repo drop (10.58\,pp) than Qwen models (5.88--6.08\,pp), reversing the function-level ranking.
We caution against generalizing from any single architecture and report these as preliminary observations requiring validation on additional architectures (e.g., CodeLlama, StarCoder).

\paragraph{Repo-level limitations.}
Multi-seed mean repo-level accuracy (79.88\%) is only 6.92\,pp above the majority baseline (72.96\%).
Fail F1 variance across seeds (std 3.87\,pp) indicates instability.
The high class imbalance (73\% pass rate) makes this a rare-event detection problem where precision-recall tradeoffs are unfavorable.
These results are preliminary; further hyperparameter tuning and data scaling may improve performance.

\paragraph{CWM token budget caveat.}
Our repo-level CWM comparison (Table~\ref{tab:cwm-token}) uses a maximum budget of 2048 tokens.
With larger budgets or better post-processing, CWM accuracy could improve further.
We frame the comparison as characterizing the efficiency--accuracy tradeoff rather than establishing a definitive accuracy ceiling for CWM.

\paragraph{Negative re-repair result.}
The predicted per-test diagnostic gain (+0.5\,pp, $p=0.72$) is not significant.
We report this honestly as a negative result: current per-test prediction quality is insufficient to translate diagnostic granularity into repair improvement.
The oracle ceiling (+4.5\,pp, $p=0.02$) suggests that improving prediction precision beyond 90\% could unlock meaningful gains---an explicit direction for future work.

%% ============================================================
%% 7  CONCLUSION
%% ============================================================
\section{Conclusion}

We presented an empirical characterization of classification and generation formulations for code execution verification.
Our findings challenge the assumption that either formulation universally dominates: classification offers efficiency advantages (1.9--8.8$\times$ speedup) and scales better to repo-level complexity, while generation captures complementary error signals.
The systematic error asymmetry---CLS conservative (75.5\% false negatives), CWM noisy (83.1\% false positives)---yields oracle ensemble ceilings of 90.63--93.1\%, motivating formulation diversity as a verification strategy.
Per-test diagnostic value is real but contingent: it helps when $T\!\ge\!2$ and prediction quality is high, but is neutral for the majority of SWE-bench instances ($T\!=\!1$, 71.2\%).

Our results are preliminary in several respects: repo-level accuracy is modest, the CLS--CWM gap is architecture-dependent, and predicted diagnostics do not yet improve repair outcomes.
We release \textsc{SWE-bench-PerTest} to enable further research on per-test verification and formulation-diverse ensemble strategies.

\bibliography{refs}
\bibliographystyle{acl_natbib}

\end{document}
